\documentclass[10pt,a4paper]{amsart}
\usepackage[margin=19mm]{geometry}
\usepackage{amsmath,amssymb,mathtools}
\usepackage[scr=rsfs,cal=euler]{mathalfa}
\usepackage{xfrac}
\usepackage[unicode,hidelinks,bookmarks=false]{hyperref}
\newcommand{\ie}{i.e.\ }
\newcommand{\cf}{cf.\ }
\newcommand{\resp}{respectively\ }
\newcommand{\Subsection}{\subsection}
\providecommand{\nobreakdash}{\nobreak\hbox{-}}
\setlength{\emergencystretch}{2em}
\multlinegap0pt
\usepackage{amssymb}
\usepackage{tensor}
\usepackage{mathtools}
\usepackage{tikz}
\newtheorem{proposition}{Proposition}
\title{\textbf{Affine Nilpotency and Engel’s Theorem for Lie Affgebras}}
\author{Tarik Anowar, Ripan Saha, Sayan Thokdar}
\date{Source: arXiv:2608.15190v1 (CC BY 4.0)}
\begin{document}
\maketitle

\noindent\emph{Source and licence.} \textbf{Affine Nilpotency and Engel’s Theorem for Lie Affgebras}. Version arXiv:2608.15190v1 (CC BY 4.0), distributed by arXiv under the Creative Commons Attribution 4.0 International licence, http://creativecommons.org/licenses/by/4.0/, as recorded in the arXiv licence metadata for this exact version and shown on the abstract page https://arxiv.org/abs/2608.15190v1. Full licence notice: \texttt{licenses/arxiv-2608.15190-CC-BY-4.0.txt}.\par\smallskip

\noindent\textit{Editorial scope.} Counted unit: the article's proposition proposition (source0.tex lines 468--470). The article's complete original proof of it is delivered with it (source0.tex lines 471--510, one \texttt{proof} environment of 40 lines). No mathematics is written, repaired or re-derived: the statement and the proof are the article's own lines, in the article's own notation, and no retained line is substituted or re-flowed.  No further source-local context is needed: every cross-reference of every retained line resolves inside the two delivered spans.  Nothing was omitted from any retained span, so the package carries no omission record.  No retained line contains a citation, so the wrapper carries no bibliography and no bibliography entry.  No retained line contains a figure environment, an image-inclusion command or drawing code, so no image is delivered and the package declares no asset dependency.  Editorial formatting only, in addition: an AMS \texttt{amsart} preamble that restates the article's own declarations and macros used by the retained lines, namely the environments \texttt{proposition} on separate counters, together with the standard packages those lines need.  The retained environments print in this document as Proposition 1 in order of appearance, whereas the article prints the same items with the numbering of its own full document.  The complete native source of the article as distributed in the arXiv e-print for this exact version is bundled unchanged as \texttt{sources/207718/source0.tex}.
\begin{proposition}
   Let $\mathfrak{a}$ be a abelian Lie affgebra with $b\in \mathfrak{a}$ such that $\{a, b\}=\{a, a\}$ for all $a\in \mathfrak{a}$. Then $b\in \mathfrak{a}$ if and only if $\lambda=\text{ad}_b$, for some $b\in \mathfrak{a}$ and $\kappa=0.$ Moreover, for any other $b'(\neq b)\in \mathfrak{a}$, $b-b'\in Z(\mathfrak{g}), \text{center of the underlying Lie algebra}.$
\end{proposition}

\begin{proof}
    Let $b\in \mathfrak{a} = \mathfrak{Z(a)}$, the center of the Lie affgebra $\mathfrak{a}$. Then, we have $\{a, b\}=\{b,a\}=\{a, a\}$, for all $a\in \mathfrak{a}$. Now, for any $a, b\in \mathfrak{a}$,
    \begin{align*}
        &\{a,b\}=[a, b]+\kappa(a)+\lambda(b-a)+s,\\
         &\{b,a\}=[b, a]+\kappa(b)+\lambda(a-b)+s,\\
         &\{a,a\}=\kappa(a)+s.
    \end{align*}
    
From the first two equality, we obtain the following,

\begin{align}
&\{a, b\}=\{b,a\},\nonumber\\
&[a, b]+\kappa(a)+\lambda(b-a)+s=[b, a]+\kappa(b)+\lambda(a-b)+s\nonumber\\
&[a, b]-[b,a]+\kappa(a-b)+2\lambda(b-a)=0.
\end{align}
From the first and third equality, we obtain the following,
\begin{align}
&[a, b]+\kappa(a)+\lambda(b-a)+s=\kappa(a)+s,\nonumber\\
&[a, b]+\lambda(b-a)=0,\nonumber\\ \label{equ13}
&[a, b]=\lambda(a-b).
\end{align}

From the second and third equality, we obtain the following,
\begin{align}
&[b, a]+\kappa(b)+\lambda(a-b)+s=\kappa(a)+s,\nonumber\\
&[b,a]+(\kappa-\lambda)(b-a)=0,\nonumber \\\label{equ14}
&[b,a]=(\kappa-\lambda)(a-b).
\end{align}
We know that in the Lie algebra $\mathfrak{g}$, by anti-commutativity, $[a, b]=-[b,a]$ for all $a, b\in \mathfrak{g}.$ Thus, using Equation $\ref{equ13}$ and $\ref{equ14}$, we get 
\begin{align}
&\lambda(a-b)=-(\kappa-\lambda)(a-b),\nonumber\\
&\kappa(a-b)=0, ~\text{for all}~ a,b \in \mathfrak{g},\nonumber\\ \label{equ15}
&\kappa(a)=\kappa(b),~\text{for all}~ a,b \in \mathfrak{g}.
\end{align}
Choosing $a=o\in \mathfrak{g},$ which is the additive identity of the Lie algebra $\mathfrak{g}$, we have $ k(b)=o$. Therefore, using Equation \ref{equ15}, we have $k(a)=0,$ for all $a\in \mathfrak{g}$. Thus, $\kappa=0.$
Using Equation $\ref{equ13}$, for all $a$, $[a, b]=\lambda(a-b)$. Again, we choose $a=o\in \mathfrak{g}$. Then, $[o, b]=o=-\lambda(b)$ and hence, $[a, b]=\lambda(a)$. Similarly, Using equation $\ref{equ14}$, for all $a$, $[b, a]=\lambda(b-a)$. Again, we choose $a=o\in \mathfrak{g}.$ Then, $[b, o]=o=\lambda(b)$ and hence $[b, a]=-\lambda(a)$. Therefore, $\text{ad}_b=\lambda.$\\
Conversely, let us assume that $\kappa=0$ and  $\lambda=\text{ad}_b$, for some $b\in \mathfrak{a}$.Then  $\{a, b\}=[a,b]+\kappa(a)+\lambda(b-a)+s=[a,b]+\text{ad}_b(b-a)+s=[a,b]+[b-a, b]+s=[a, b]+[b, b]-[a, b]+s=s$, $\{b, a\}=[b,a]+\kappa(b)+\lambda(a-b)+s=[b,a]+\lambda(a-b)+s=[b,a]+\text{ad}_b(a-b)+s=[b,a]+[a-b, b]+s=[b,a]+[a, b]-[b,b]+s=s$ and $\{a, a\}=s,$ for all $a\in \mathfrak{a}.$ Then for all $a\in \mathfrak{a}$, $\{a, b\}=\{b,a\}=\{a, a\}.$\\
Let $b'(\neq b) \in \mathfrak{Z(a)}$, then by above observation, $\text{ad}_b(a)=\text{ad}_{b'}(a)=\lambda(a)$. It follows that $[a, b]=[a, b']$ and then $[a, b-b']=0, \forall a \in\mathfrak{a}$. Thus, $b-b'\in Z(\mathfrak{g}).$
    
\end{proof}

\end{document}
